\section{Arithmetic positivity at the bounded scale}
\label{sec:arithmetic}

We continue the minimal-counterexample argument with a low factor.  Let
\(Z_i=\prod_{\ell=0}^s Z_{i\ell}\) be the subproducts supplied by
Proposition~\ref{prop:minimal-products}, and write
\[
 d_\ell=\dim Z_{i\ell},\qquad d=\sum_{\ell=0}^s d_\ell,
 \qquad r=\dim\cA+1.
\]
The dimensions are constant after extraction.
The goal is an arithmetic lower bound for these varying products. Its
proof compares two measures on the fixed low factor: a measure from
separate outputs and one from consecutive differences. The model is the
metric-variation method of Szpiro, Ullmo and Zhang
\cite[Theorem~3.1 and its proof]{SzpiroUllmoZhang1997}, in the arithmetic
bigness form developed by Yuan \cite[Section~3.2]{Yuan2008}. We give the
uniform estimates here because both the high factors and their fields
of definition vary.  Lemma~\ref{lemma:fixed-low-point}
has disposed of the case \(d_0=0\).  Thus
\(Z_{i0}=Z_0\) is fixed and \(d_0>0\).  Each product is geometrically
integral over a number field \(k_i\), the degrees \([k_i:\mathbb Q]\)
are bounded, and the fields contain all fixed ambient data.  Enlarge them
by \(\mathbb Q(\sqrt{-1})\), preserving this bound.  They carry the
embedding in \(\mathbb C\) fixed at the start of the proof.  The degrees
of all factors are bounded, and each \(Z_{i\ell}\) has a dense set of
points whose projective heights are at most \(C H_{i\ell}\), with one
constant \(C\).  Recall that \(H_{i0}=1\).  These bounds concern the
auxiliary varieties and their defining fields; the marked points may have
arbitrarily large degree.

\subsection{Scaled bundles and their geometric degrees}

Retain the symmetric very ample line bundles \(L_\ell\) on \(A_\ell\)
and fix a symmetric very ample line bundle \(L\) on \(\cA\), all with their
canonical adelic metrics.  We denote the resulting
nef adelic bundles by \(\bar L_\ell\) and \(\bar L\).  All arithmetic
intersection numbers and heights below are absolute: intersection numbers
over \(k_i\) are divided by \([k_i:\mathbb Q]\).

After another extraction, the homology class of every high factor
\(Z_{i\ell}(\mathbb C)\) is constant.  Indeed, integrals of fixed smooth
forms are bounded in terms of the projective degree, so these integral
homology classes lie in a bounded subset of a lattice.  Choose positive
integers \(M_i\to\infty\) and \(b_{i\ell}\) such that
\[
 b_{i0}=M_i,\qquad
 b_{i\ell}\sim\frac{M_i}{\sqrt{H_{i\ell}}},\qquad
 \min_\ell b_{i\ell}\longrightarrow\infty.
\]
Because \(M_i\) is unrestricted, lattice approximation in the fixed
homomorphism groups gives homomorphisms \(F_i:\cA\to\cA\) with
\begin{equation}\label{eq:scaled-homomorphisms}
 \varphi_i:=F_i\operatorname{diag}(b_{i\ell}^{-1})
 \longrightarrow\Phi.
\end{equation}
The maps \(F_i\) may be chosen to be isogenies: first approximate
\(\Phi\) by invertible real endomorphisms, then choose the scaling
integers sufficiently large that lattice approximation preserves
invertibility.  Set
\[
 \bar E_i=F_i^*\bar L,\qquad
 B_i=\prod_{\ell=0}^s b_{i\ell}^{2d_\ell}.
\]
The normalized marked tuples are bounded, and their images under \(\Phi\)
tend to zero.  Together with \eqref{eq:scaled-homomorphisms}, this gives
\begin{equation}\label{eq:small-marked-height}
 h_{\bar E_i}(y_i)=o(M_i^2).
\end{equation}
More explicitly, the vector with coordinates
\(b_{i\ell}y_{i\ell}/M_i\) differs by a vector tending to zero from
\((y_{i\ell}/\sqrt{H_{i\ell}})_\ell\); apply \(\varphi_i\) and use
the norm comparison for real homomorphisms.

On \(\cA^r\), let \(\bar E_{i,1}\) be the pullback of the external
product of \(r\) copies of \(\bar L\) under
\[
 (x_1,\ldots,x_r)\longmapsto(F_i(x_1),\ldots,F_i(x_r)),
\]
and let \(\bar E_{i,2}\) be the analogous pullback for the \(r-1\)
consecutive differences \(F_i(x_{a+1})-F_i(x_a)\).  Both are nef.
The tangent positivity in Proposition~\ref{prop:tangent}, the block
scaling identity for top forms, and the fixed cycle classes give
\begin{equation}\label{eq:arithmetic-geometric-degrees}
 (E_i^d\cdot Z_i)\asymp B_i,
 \qquad
 (E_{i,j}^{rd}\cdot Z_i^r)\asymp B_i^r\quad(j=1,2).
\end{equation}
The lower bounds here are uniform.  After division by the displayed
scaling factors, the degrees converge to the positive integrals for
\(\Phi\) and its consecutive-difference map on the fixed cycle classes.

\subsection{Two measures and arithmetic positivity}

Although the marked points have \(\bar E_i\)-height \(o(M_i^2)\), the
following proposition shows that the normalized arithmetic intersections
of the whole varieties have a positive lower bound.  We will compare the
curvature measures of the separate-output and difference bundles on the
fixed low product \(Z_0^r\).

\begin{proposition}\label{prop:arithmetic-positive}
Let \(Z_i=\prod_{\ell=0}^s Z_{i\ell}\) be the bounded-degree,
bounded-field-degree subproducts above, with fixed positive-dimensional
low factor \(Z_0\), dense factorwise height bounds
\(h_{L_\ell}=O(H_{i\ell})\), and fixed high cycle classes.  Suppose the
two tangent integrals for \(\Phi\) in
Proposition~\ref{prop:tangent} are positive.  For the homomorphisms and
scalings in \eqref{eq:scaled-homomorphisms}, one has
\begin{equation}\label{eq:arithmetic-positive}
 \liminf_{i\to\infty}
 \frac{(\bar E_i^{d+1}\cdot Z_i)}{M_i^2B_i}>0.
\end{equation}
\end{proposition}

\begin{proof}
We first compute two distinct limiting measures geometrically.  We then
show that a vanishing subsequence of the arithmetic intersection
quotients would force those same measures to agree.

\emph{The two measures and their geometric limits.}

Put \(X_i=Z_i^r\), \(u=rd\), and let
\(p_0:\cA^r\to A_0^r\) denote projection to the low coordinates.
For \(j=1,2\), define a probability measure on \(A_0^r(\mathbb C)\) by
\[
 \mu_{i,j}
 =\frac{(p_0)_*\big(c_1(\bar E_{i,j})^u\big|_{X_i}\big)}
 {(E_{i,j}^u\cdot X_i)}.
\]
Curvatures in this formula are taken at the chosen complex embedding.
Both measures are supported on the fixed low product \(Z_0^r\).
The geometric comparison follows the product-versus-pullback argument
of the Bogomolov method \cite[Section~4]{Ullmo1998}; the
consecutive-difference construction is explained in
\cite[Section~3]{Zhang1998ICM}, following Faltings.

Let \(\omega=c_1(\bar L)\) at the fixed complex embedding.  For a real
homomorphism \(g\), let \(\beta_1(g)\) be the invariant form on
\(\cA^r\) obtained by summing the pullbacks of \(\omega\) under \(g\)
in the \(r\) separate copies.  Let \(\beta_2(g)\) be the sum of the
pullbacks under the consecutive differences of those outputs.  These
definitions use only linear maps on Lie algebras and therefore make sense
also for \(g=\Phi\).

On a product tangent space, block scaling multiplies the top form by
\(B_i^r\).  For any smooth real function \(f\) on
\(A_0^r(\mathbb C)\), multiplication by \(p_0^*f\) does not alter
that pointwise identity.  It follows that
\begin{equation}\label{eq:remove-scales-with-test}
 \mu_{i,j}(f)
 =\frac{\displaystyle\int_{Z_i^r}p_0^*f\,
                            \beta_j(\varphi_i)^u}
        {\displaystyle\int_{Z_i^r}\beta_j(\varphi_i)^u}.
\end{equation}
No multiplication map on points or change of the test function is being
used in this identity.

Write \(Q_i=\prod_{\ell=1}^s Z_{i\ell}\), allowing a point when
there are no high factors.  Integration of an invariant form along
\(Q_i^r\) gives an invariant form on \(A_0^r\): in invariant
coordinates its coefficients are integrals of fixed invariant forms
over \(Q_i^r\), hence depend only on its homology class.  Choose one
of the products \(Q_i\) as a fixed representative \(Q\).  Integrating
\(\beta_j(g)^u\) over \(Q^r\) produces a smooth invariant form
\(\gamma_j(g)\) of bidegree \((rd_0,rd_0)\) on \(A_0^r\).
Its coefficients are polynomial in those of \(g\).  Thus
\(\gamma_j(\varphi_i)\to\gamma_j(\Phi)\) smoothly, and
\eqref{eq:remove-scales-with-test} shows that \(\mu_{i,j}\) converges
to the probability measure \(\mu_j\) given by
\[
 \mu_j(f)
 =\frac{\displaystyle\int_{Z_0^r}f\,\gamma_j(\Phi)}
        {\displaystyle\int_{Z_0^r}\gamma_j(\Phi)}.
\]
Both denominators are positive by the tangent positivity already used
in \eqref{eq:arithmetic-geometric-degrees}.  Singularities cause no
problem in these integrations: one may integrate over resolutions,
and the resulting densities on \((Z_0^{\mathrm{sm}})^r\) are smooth.

The first tangent integral supplies smooth points \(z\in Z_0\) and
\(q\in Q\) such that \(\Phi\) is injective on
\(T_zZ_0\oplus T_qQ\), where
the tangent spaces are translated into the ambient Lie algebras.
At the diagonal low tuple \((z,\ldots,z)\), the density
\(\gamma_1(\Phi)|_{Z_0^r}\) is strictly positive.  Indeed, at
\((q,\ldots,q)\) the separate-output form is positive on the full
product tangent space, and it remains positive on a nonempty open set
of high tuples, whose integration gives a positive density.

At this same low tuple the density
\(\gamma_2(\Phi)|_{Z_0^r}\) is zero.  Choose a nonzero tangent vector
\(v\in T_zZ_0\), possible because \(d_0>0\).
For every smooth high tuple, the tangent vector which has low part
\((v,\ldots,v)\) and zero high parts is killed by every consecutive
difference.  The top power of the semipositive form
\(\beta_2(\Phi)\) therefore vanishes on that full product tangent
space; its integral over the high factors vanishes as well.

After their positive normalizations, the two smooth densities still
differ at \((z,\ldots,z)\).  Their difference has a fixed positive
sign on a sufficiently small neighborhood in the smooth locus.  Choose
a nonnegative smooth function \(f_*\) on \(A_0^r(\mathbb C)\),
supported in a coordinate neighborhood of this point and positive on a
smaller neighborhood.  We have proved
\begin{equation}\label{eq:low-measure-gap}
 \lim_i\bigl(\mu_{i,1}(f_*)-\mu_{i,2}(f_*)\bigr)>0.
\end{equation}
Figure~\ref{fig:low-measures} records this comparison on the low product.
\begin{figure}[htbp]
\centering
\begin{tikzpicture}[
  box/.style={draw,align=center,text width=5.2cm,
              inner sep=6pt,font=\small},
  arrow/.style={-{Stealth[length=2mm]},thick},
  note/.style={font=\footnotesize,align=center}]
\node[align=center,font=\small] (source) at (0,0)
  {Common source $X_i=Z_0^r\times Q_i^r$};
\node[box] (one) at (-3,-1.1)
  {Separate-output curvature\\$c_1(\bar E_{i,1})^u|_{X_i}$};
\node[box] (two) at (3,-1.1)
  {Consecutive-difference curvature\\$c_1(\bar E_{i,2})^u|_{X_i}$};
\draw[arrow] (source.south) -- (one.north);
\draw[arrow] (source.south) -- (two.north);
\node[box] (lowone) at (-3,-3.1)
  {$\mu_{i,1}\longrightarrow\mu_1$ on $Z_0^r$\\
   positive density at $(z,\ldots,z)$};
\node[box] (lowtwo) at (3,-3.1)
  {$\mu_{i,2}\longrightarrow\mu_2$ on $Z_0^r$\\
   zero density at $(z,\ldots,z)$};
\draw[arrow] (one) -- node[note,fill=white,inner sep=3pt]
  {normalize; apply $(p_0)_*$} (lowone);
\draw[arrow] (two) -- node[note,fill=white,inner sep=3pt]
  {normalize; apply $(p_0)_*$} (lowtwo);
\node at (0,-3.1) {$\ne$};
\end{tikzpicture}
\caption{The two curvature measures on the changing product $X_i$ are
normalized by their geometric degrees and projected to the fixed low
product $Z_0^r$. Their limits differ: a common nonzero low tangent vector
is killed by consecutive differences. If the arithmetic intersection
quotient $(\bar E_i^{d+1}\cdot Z_i)/(M_i^2B_i)$ tended to zero, uniform
metric variation would force these limits to agree. The smooth density
comparison gives the test function in \eqref{eq:low-measure-gap}.}
\label{fig:low-measures}
\end{figure}


\emph{Small points from a vanishing arithmetic intersection.}
We now suppose, contrary to the proposition, that the nonnegative
quotient in \eqref{eq:arithmetic-positive} tends to zero along a
subsequence.  We recall the two arithmetic inequalities used here and
in the small-section construction below.  For a nef adelic bundle
\(\bar D\) on a projective variety \(X\) of dimension \(n\), with
\(D\) geometrically big, the fundamental inequalities read
\begin{equation}\label{eq:fundamental-inequalities}
 \frac{e_1(X,\bar D)}{n+1}
 \le \frac{(\bar D^{n+1}\cdot X)}{(n+1)(D^n\cdot X)}
 \le e_1(X,\bar D).
\end{equation}
Here \(e_1\) is the essential minimum of the point height.  We use the
nef, geometrically big form of these inequalities
\cite[Theorem~5.3.3]{YuanZhang2021}; the ample projective case is
\cite[Theorem~1.10]{Zhang1995}.  We also use the arithmetic Siu inequality
\begin{equation}\label{eq:arithmetic-siu}
 \widehat{\operatorname{vol}}(\bar D-\bar N)
 \ge (\bar D^{n+1}\cdot X)
       -(n+1)(\bar D^n\bar N\cdot X)
\end{equation}
for nef adelic bundles, with the same absolute normalization
\cite[Theorem~A.5.1(2)]{YuanZhang2021}.
These statements apply after a projective resolution; pullback preserves
the intersection numbers and essential minima used here. Resolution in
characteristic zero can be performed over the original field \(k_i\),
by smooth-center blowups of a smooth projective ambient variety, with
centers disjoint from the original smooth locus of \(X\)
\cite[Theorems~1.0.2--1.0.3 and Section~5.7]{Wlodarczyk2005}.
Thus using a resolution
requires no extension of the bounded-degree field \(k_i\).  A bundle of
positive arithmetic volume has a nonzero section of some positive power
with supremum norm at most one at every place.  Its point height is
therefore nonnegative outside a proper closed subset.

By \eqref{eq:fundamental-inequalities} and
\eqref{eq:arithmetic-geometric-degrees}, each \(Z_i\) has a dense
set of points of \(\bar E_i\)-height at most \(\epsilon_iM_i^2\),
where \(\epsilon_i\to0\).  Products of these sets give a dense subset
\(\mathcal S_i\subset X_i(\Qbar)\) such that
\begin{equation}\label{eq:simultaneous-small-height}
 h_{\bar E_{i,1}}(x),\ h_{\bar E_{i,2}}(x)
 \le C_r\epsilon_iM_i^2\qquad(x\in\mathcal S_i).
\end{equation}
For the difference bundle this follows from
\(\|a-b\|^2\le2\|a\|^2+2\|b\|^2\).
We will use these common dense sets to show that, for every smooth real
function \(f\) on \(A_0^r(\mathbb C)\),
\begin{equation}\label{eq:low-measure-comparison}
 \mu_{i,1}(f)-\mu_{i,2}(f)\longrightarrow0.
\end{equation}
Applied to \(f_*\), this will contradict \eqref{eq:low-measure-gap}.
The required metric variation must be uniform in both the varieties and
their defining fields.

\emph{Uniform bounds for metric variation.}
Let \(v_i\) be the complex place of \(k_i\) determined by the chosen
embedding, and put \(\lambda_i=2/[k_i:\mathbb Q]\).  Thus
\(\lambda_i\) is bounded below by a positive constant.  Write
\(\overline{\mathcal O}_{v_i}(f)\) for the trivial bundle whose negative
logarithmic norm is \(p_0^*f\) at this embedding, its conjugate at the
conjugate embedding, and zero at all other places.  Write
\(\bar{\mathbf 1}\) for the trivial bundle with negative logarithmic
norm one at every archimedean place and zero at the finite places.
Its absolute point height is one.

Fix \(\eta>0\).  For either \(j\), abbreviate
\(\bar D_i=\bar E_{i,j}|_{X_i}\), \(\mu_i=\mu_{i,j}\), and set
\[
 a_i=\lambda_i\mu_i(f)-\eta,\qquad
 \bar U_i=\overline{\mathcal O}_{v_i}(f)-a_i\bar{\mathbf 1}.
\]
In particular \(|a_i|\le\|f\|_\infty+\eta\).  There are nef bundles
\(\bar U_{i,1},\bar U_{i,2}\), pulled back from the low ambient product,
such that
\begin{equation}\label{eq:uniform-nef-decomposition}
 \bar U_i=\bar U_{i,1}-\bar U_{i,2},\qquad
 \bar U_{i,2}=C\bar H+K\bar{\mathbf 1},
\end{equation}
where \(\bar H\) is a fixed canonical symmetric ample bundle on
\(A_0^r\), and \(C,K\) are fixed positive constants.  To obtain this,
choose an integer \(C\) so that
\(C c_1(\bar H)+dd^cf\) is semipositive, and then choose
\(K>\|f\|_\infty+\sup_i|a_i|\).

We spell out why this is nef in the model sense required by
\eqref{eq:arithmetic-siu}.  Start with a relatively ample projective model
of a power of \(H\), give it the invariant archimedean curvature, and
add a constant so that it is nef.  Normalized pullbacks by multiplication,
on graph models, converge to \(\bar H\); their archimedean curvatures
are all the same invariant form.  Adding the metric in
\eqref{eq:uniform-nef-decomposition} leaves every vertical curve degree
unchanged. At infinity the new curvature is semipositive by the choice
of \(C\). The added negative-log-norm function is nonnegative at each
archimedean place by the choice of \(K\), so the arithmetic degree of
every horizontal integral curve can only increase. This verifies the
nef criterion for hermitian models: semipositive curvature and
nonnegative degrees on all integral arithmetic curves. The perturbed
nef models converge to \(\bar U_{i,1}\).
The construction also applies to \(\bar U_{i,2}\), and its constants
are independent of \(i\) and of the field extension.

Every arithmetic mixed intersection of total order \(u+1\) formed from
\[
 \bar D_i,\qquad M_i^2\bar U_{i,1},\qquad M_i^2\bar U_{i,2}
\]
is bounded by
\begin{equation}\label{eq:uniform-mixed-intersections}
 O(M_i^2B_i^r),
\end{equation}
with one constant for all these finitely many types of intersections.
For a direct proof, let \(p_{a\ell}\) be the coordinate projection from
\(\cA^r\) to factor \(\ell\) in copy \(a\), and put
\[
 \bar T_i=\sum_{a=1}^r\sum_{\ell=0}^s
 b_{i\ell}^2p_{a\ell}^*\bar L_\ell,
 \qquad
 \bar S_i=\bar D_i+M_i^2\bar U_{i,1}
                 +M_i^2\bar U_{i,2}+\bar T_i
\]
on \(X_i\).  This is nef and geometrically ample.  Remove the factors
\(b_{i\ell}\) on each product tangent space.  The remaining forms
for \(D_i\) are bounded by \eqref{eq:scaled-homomorphisms}; those for
\(M_i^2U_{i,a}\) are bounded because their geometric parts use only
the low coordinates, where \(b_{i0}=M_i\).  Consequently
\[
 (S_i^u\cdot X_i)=O(B_i^r).
\]
On the products of the dense factorwise sets of height
\(O(H_{i\ell})\), all the point heights just displayed are
\(O(M_i^2)\).  For \(D_i\), use boundedness of \(\varphi_i\) to get
\[
 h_{\bar D_i}(x)\ll
 \sum_{a,\ell} b_{i\ell}^2
                 h_{\bar L_\ell}(x_{a\ell})=O(M_i^2).
\]
The same estimate holds for \(\bar T_i\).  The heights of
\(\bar U_{i,1},\bar U_{i,2}\) are bounded on these sets because the
low point heights and the metric shifts are bounded.  Thus
\(e_1(X_i,\bar S_i)=O(M_i^2)\), and the upper inequality in
\eqref{eq:fundamental-inequalities} yields
\((\bar S_i^{u+1}\cdot X_i)=O(M_i^2B_i^r)\).
Mixed intersections of nef adelic bundles are nonnegative.  Each term
in \eqref{eq:uniform-mixed-intersections} occurs with a positive
coefficient in this last self-intersection, which proves the assertion.
Absolute normalization has introduced no field-dependent constant.

\emph{Comparison against a low-factor test function.}
The metric intersection formula gives
\[
 (\bar D_i^u\bar U_i\cdot X_i)
 =(D_i^u\cdot X_i)\big(\lambda_i\mu_i(f)-a_i\big)
 =\eta(D_i^u\cdot X_i).
\]
Apply \eqref{eq:arithmetic-siu} to the two nef bundles
\(\bar D_i+tM_i^2\bar U_{i,1}\) and
\(tM_i^2\bar U_{i,2}\), where \(t>0\) is rational.  Expansion and
\eqref{eq:uniform-mixed-intersections} give
\begin{align*}
 \widehat{\operatorname{vol}}(\bar D_i+tM_i^2\bar U_i)
 \ &\ge (\bar D_i^{u+1}\cdot X_i)
 +(u+1)tM_i^2\eta(D_i^u\cdot X_i)
 -O(t^2M_i^2B_i^r).
\end{align*}
The first term is nonnegative.  By
\eqref{eq:arithmetic-geometric-degrees}, a sufficiently small fixed
\(t>0\), depending on \(f,\eta\) but not on \(i\), makes the right
side positive for all sufficiently large \(i\).  Hence the perturbed
height is nonnegative outside a proper closed subset of \(X_i\).

Repeat this argument for both \(j\) and for \(f\) and \(-f\).
Choose \(x_i\in\mathcal S_i\) outside the four resulting proper
closed subsets.  Let \(\operatorname{av}_{i,x_i}(f)\) be the average
of \(f\) on the low projections of the conjugates of \(x_i\) over
\(k_i\) at the chosen embedding.  Then
\[
 h_{\overline{\mathcal O}_{v_i}(f)}(x_i)
 =\lambda_i\operatorname{av}_{i,x_i}(f).
\]
Divide the perturbed height inequalities by \(tM_i^2\), and use
\eqref{eq:simultaneous-small-height}.  For \(j=1,2\) this yields
\[
 \left|\lambda_i\operatorname{av}_{i,x_i}(f)
             -\lambda_i\mu_{i,j}(f)\right|
 \le\eta+o(1).
\]
Since \(\inf_i\lambda_i>0\) and \(\eta\) is arbitrary, this proves
\eqref{eq:low-measure-comparison}.  This is the only point in this
metric-variation argument that needs the bounded degrees of the fields
\(k_i\); the degrees of the points \(x_i\) are unrestricted.

Taking \(f=f_*\) in \eqref{eq:low-measure-comparison} contradicts
\eqref{eq:low-measure-gap}.  This proves
\eqref{eq:arithmetic-positive}.
\end{proof}

\subsection{Sections with a uniform exponential bound}

The arithmetic intersection lower bound permits subtracting a constant
adelic bundle while keeping positive arithmetic volume.  Returning to
the original metric gives the following small sections over the fields
\(k_i\).

\begin{lemma}\label{lemma:small-sections}
For the products \(Z_i\), bundles \(\bar E_i\), and scalings above,
there is a constant \(c>0\) such that, for all sufficiently large
\(i\), a projective resolution \(\pi_i:\widetilde Z_i\to Z_i\)
over \(k_i\) carries a nonzero section defined over \(k_i\),
\[
 s_i\in H^0(\widetilde Z_i,\pi_i^*E_i^{\otimes n_i})
 \quad\text{for some integer }n_i>0
\]
whose adelic supremum norms satisfy
\begin{equation}\label{eq:small-section-norms}
 \|s_i\|_{v,\mathrm{sup}}\le
 \begin{cases}
 \exp(-cn_iM_i^2),&v\mid\infty,\\
 1,&v\nmid\infty.
 \end{cases}
\end{equation}
The resolution may be chosen to be an isomorphism over the smooth locus.
\end{lemma}

\begin{proof}
By Proposition~\ref{prop:arithmetic-positive} and
\eqref{eq:arithmetic-geometric-degrees}, there are constants
\(\delta,C>0\) with
\[
 (\bar E_i^{d+1}\cdot Z_i)\ge\delta M_i^2B_i,
 \qquad (E_i^d\cdot Z_i)\le CB_i.
\]
Choose \(0<c<\delta/((d+1)C)\).  The constant bundle
\(cM_i^2\bar{\mathbf 1}\) is nef, and
\eqref{eq:arithmetic-siu} shows that
\(\bar E_i-cM_i^2\bar{\mathbf 1}\) has positive arithmetic volume.
After pulling back to the resolution, take a nonzero section of a
positive power with norms at most one for this shifted metric.  Restoring
the original metric gives exactly \eqref{eq:small-section-norms}.
\end{proof}

The exponent \(n_i\) need not be bounded.  Its role is homogeneous:
the logarithmic norm bound is proportional to \(n_iM_i^2\), and the
weighted order in the next section is normalized by \(n_i b_{i\ell}^2\).
The fixed positive constant \(c\), together with
\(h_{\bar E_i}(y_i)=o(M_i^2)\), will force a positive weighted index
at the marked point.  Enlarging the field to calculate at that point
will not change the field of definition of the section.
